% Options for packages loaded elsewhere
\PassOptionsToPackage{unicode}{hyperref}
\PassOptionsToPackage{hyphens}{url}
%
\documentclass[
]{article}
\usepackage{amsmath,amssymb}
\usepackage{iftex}
\ifPDFTeX
  \usepackage[T1]{fontenc}
  \usepackage[utf8]{inputenc}
  \usepackage{textcomp} % provide euro and other symbols
\else % if luatex or xetex
  \usepackage{unicode-math} % this also loads fontspec
  \defaultfontfeatures{Scale=MatchLowercase}
  \defaultfontfeatures[\rmfamily]{Ligatures=TeX,Scale=1}
\fi
\usepackage{lmodern}
\ifPDFTeX\else
  % xetex/luatex font selection
\fi
% Use upquote if available, for straight quotes in verbatim environments
\IfFileExists{upquote.sty}{\usepackage{upquote}}{}
\IfFileExists{microtype.sty}{% use microtype if available
  \usepackage[]{microtype}
  \UseMicrotypeSet[protrusion]{basicmath} % disable protrusion for tt fonts
}{}
\makeatletter
\@ifundefined{KOMAClassName}{% if non-KOMA class
  \IfFileExists{parskip.sty}{%
    \usepackage{parskip}
  }{% else
    \setlength{\parindent}{0pt}
    \setlength{\parskip}{6pt plus 2pt minus 1pt}}
}{% if KOMA class
  \KOMAoptions{parskip=half}}
\makeatother
\usepackage{xcolor}
\usepackage[margin=1in]{geometry}
\usepackage{longtable,booktabs,array}
\usepackage{calc} % for calculating minipage widths
% Correct order of tables after \paragraph or \subparagraph
\usepackage{etoolbox}
\makeatletter
\patchcmd\longtable{\par}{\if@noskipsec\mbox{}\fi\par}{}{}
\makeatother
% Allow footnotes in longtable head/foot
\IfFileExists{footnotehyper.sty}{\usepackage{footnotehyper}}{\usepackage{footnote}}
\makesavenoteenv{longtable}
\usepackage{graphicx}
\makeatletter
\newsavebox\pandoc@box
\newcommand*\pandocbounded[1]{% scales image to fit in text height/width
  \sbox\pandoc@box{#1}%
  \Gscale@div\@tempa{\textheight}{\dimexpr\ht\pandoc@box+\dp\pandoc@box\relax}%
  \Gscale@div\@tempb{\linewidth}{\wd\pandoc@box}%
  \ifdim\@tempb\p@<\@tempa\p@\let\@tempa\@tempb\fi% select the smaller of both
  \ifdim\@tempa\p@<\p@\scalebox{\@tempa}{\usebox\pandoc@box}%
  \else\usebox{\pandoc@box}%
  \fi%
}
% Set default figure placement to htbp
\def\fps@figure{htbp}
\makeatother
\setlength{\emergencystretch}{3em} % prevent overfull lines
\providecommand{\tightlist}{%
  \setlength{\itemsep}{0pt}\setlength{\parskip}{0pt}}
\setcounter{secnumdepth}{-\maxdimen} % remove section numbering
\usepackage{bookmark}
\IfFileExists{xurl.sty}{\usepackage{xurl}}{} % add URL line breaks if available
\urlstyle{same}
\hypersetup{
  hidelinks,
  pdfcreator={LaTeX via pandoc}}

\author{}
\date{\vspace{-2.5em}}

\begin{document}

\section{RVAT NLQ Interface}\label{rvat-nlq-interface}

Version: 0.3\\
Date: 2026.09.04

A Natural Language Query (NLQ) interface for exploring RVAT genetic
variant data using locally hosted Large Language Models (LLMs).

================

\section{Introduction}\label{introduction}

\subsection{RVAT Natural Language Querey (NLQ)
Interface}\label{rvat-natural-language-querey-nlq-interface}

\subsubsection{Project Objective}\label{project-objective}

This project evaluates whether Large Language Models (LLMs) can
effectively translate natural language questions into executable SQL
queries for scientific analysis and research of datasets against the
RVAT database.

\textbf{RVAT} (Rare Variant Association Toolkit) was developed by
\href{https://kennalab.github.io/rvat/articles/basics.html}{Kevin Kenna
and Paul Hop} and is an R package and command-line toolkit designed for
large-scale genetic association testing using rare genomic variants.

The RVAT NLQ Interface allows users to interact with a database using
natural language rather than writing SQL manually.

A user's question is translated into SQL by an Ollama-hosted LLM,
executed against the target database, and the results are returned
through a Shiny-based web interface.

\subsubsection{Research Question}\label{research-question}

The primary objective is to assess:

To what extent can a locally hosted Large Language Model (LLM),
integrated within an R environment, translate natural language queries
into accurate SQL queries and use them to retrieve relevant data from an
ALS-specific RVAT database, as evaluated using predefined benchmark
questions?

The project establishes three core technical objectives:

• be used as an evaluation framework for Natural Language → SQL
generation;\\
• serve as a reference implementation of a local, privacy-preserving
LLM-driven query interface;\\
• demonstrate reproducible R/Shiny application that can be deployed in
research or server environments.\\

\begin{center}\rule{0.5\linewidth}{0.5pt}\end{center}

================

\section{Architecture Overview}\label{architecture-overview}

The application consists of six main components:

\begin{longtable}[]{@{}ll@{}}
\toprule\noalign{}
Script & Purpose \\
\midrule\noalign{}
\endhead
\bottomrule\noalign{}
\endlastfoot
\texttt{01\_db\_connection.R} & Database connection management \\
\texttt{02\_ollama\_config.R} & Ollama and LLM configuration \\
\texttt{03\_query\_execution.R} & Natural language → SQL execution
logic \\
\texttt{04\_logging\_pipeline.R} & Query logging and metadata capture \\
\texttt{05\_shiny\_ui.R} & User interface \\
\texttt{06\_shiny\_server.R} & Application server logic \\
\end{longtable}

These components are orchestrated through: rvat\_nlq\_app.R

This is the \textbf{only script that end users need to run}. All other
scripts are loaded automatically.
------------------------------------------------------------------------

================

\section{System Requirements}\label{system-requirements}

\subsection{Reference Environment}\label{reference-environment}

The application was developed and validated in the following server
environment: - Operating system : Linux(x86\_64) - R Version : 4.5.0
(2025-04-11) - CPU : 32 cores - Memory : 64GB RAM - Available memory :
\textasciitilde{} 49GB

\subsection{Hardware and Software
Specifications}\label{hardware-and-software-specifications}

\textbf{Server Environment (Recommended)}

For full reproducibility:

• Linux Server (x86\_64) • 32 CPU cores • 64 GB RAM • RStudio 4.5.0
(2025-04-11) • Ollama installed (0.21.0) • Ollamar installed (1.2.2) •
One or more supported LLM's

This configuration was used throughout development and validation.

\textbf{Local/Laptop Environment} Experimental trial: Local execution is
supported but depends on available hardware.

Smaller models generally perform well:

mistral:latest qwen2.5-coder:latest

Large models may experience reduced performance or even freeze:

qwen3:8b sqlcoder:latest

Depending on available CPU and RAM these models may:

\begin{itemize}
\tightlist
\item
  require several minutes to respond
\item
  consume substantial resources
\item
  produce unstable behaviour during long-running requests
\end{itemize}

This behaviour is related to hardware constraints rather than
application logic.

Laptop Reference Specifications: • Dell Latitude 5220 • Windows 11 Pro
(x64) • CPU: 4 cores • 8GB RAM • Disk space to download Ollama/LLM's
locally

** Windows \& Bioconductor Note ** RStudio and R do not natively require
Bioconductor on Windows. However, core genomic dependencies used in this
project (such as \texttt{SummarizedExperiment}, required upstream by
\texttt{rvat}) frequently fail to resolve automatically on Windows via
standard \texttt{install.packages()}.

While macOS and Linux typically handle these upstream Bioconductor
dependencies dynamically, Windows environments require explicit
initialization of \texttt{BiocManager} to register the necessary
repositories and build binary packages smoothly.

** Setup Script for Laptop Users ** To streamline setup and resolve
cross-platform dependency issues (especially on Windows), an automated
setup script is included in the project root:
\texttt{/scripts/00\_renv\_setup.R}.

This script automatically: * Verifies and installs \texttt{renv} and
\texttt{BiocManager} (ensuring correct Bioconductor repository mapping).
* Restores the exact project dependencies via \texttt{renv::restore()}.
* Conducts post-installation checks to ensure core packages
(\texttt{shiny}, \texttt{dplyr}, \texttt{DBI}, \texttt{RSQLite},
\texttt{rvat}) are ready.

\begin{center}\rule{0.5\linewidth}{0.5pt}\end{center}

\section{Getting Started}\label{getting-started}

Before starting, ensure to review `Hardware and Software Specifications'
above to your chosen system:

\subsection{System Requirements}\label{system-requirements-1}

\begin{itemize}
\tightlist
\item
  R (≥ 4.x recommended)
\item
  RStudio / Posit
\item
  Git
\item
  Internet connection (for package and model download)
\end{itemize}

\subsection{Required Tools}\label{required-tools}

\begin{itemize}
\item
  \textbf{Ollama (mandatory)}\\
  👉 \url{https://ollama.com/}
\item
  At least one LLM model with tool support\\
  Recommended:
\end{itemize}

```bash ollama pull qwen3:8b

\end{document}
